\section{The exact restriction theorem}


\begin{theorem}[Consecutive quartet domination]
\label{thm:reduction}
Choose any cheapest adjacent pair $e$ of cost $q=\OPT_{N-1}$ and any
contiguous flat optimum $P$ at capacity $N-2$, of cost $F=\OPT_{N-2}$.
One of the following conclusions holds.
\begin{enumerate}[label=\textup{(\roman*)}]
\item There is a hierarchy with both levels flat optimal, so $R_N=1$.
\item The only nonsingleton cells of $P$ are two pairs $AB$ and $CD$ on
four consecutive original states $A<B<C<D$, and $e=BC$. For the source
restricted to $S=\{A,B,C,D\}$, keeping its original weights, its flat
values satisfy
\[
 \OPT_3(S)=q,\qquad \OPT_2(S)=F=D(AB)+D(CD),
\]
and its unrestricted deterministic $(2,3)$ hierarchy price satisfies
\[
 R_N\le R_4(S).
\]
\end{enumerate}
In particular, if $R_N>1$, conclusion \textup{(ii)} holds for every such
choice of $e$ and $P$. When $N=3$, necessarily $R_N=1$.
\end{theorem}

\begin{proof}
The partition whose sole nonsingleton cell is $e$ is flat optimal at the
fine capacity. We first check all possible forms of $P$.

Suppose the sole nonsingleton cell of $P$ is a triple $T=ABC$ of consecutive
states. If both endpoints of $e$ belong to $T$, the fine partition just
described refines $P$, proving (i). If $e$ is disjoint from $T$, take
$f=AB$ and $h=BC$. Then $e$ and $f$ are disjoint, and
\[
 D(e)+D(f)\le D(h)+D(f)\le D(T)=F,
\]
because $e$ is globally cheapest among adjacent pairs and
Lemma~\ref{lem:triple} applies. The partition with nonsingleton cells
$e,f$ has $N-2$ cells and cost at most $F$, hence exactly $F$. It is a
flat optimum compatible with the chosen fine partition. If $e$ meets
$T$ in exactly one state, adjacency of $e$ and consecutivity of $T$
force that intersection to be an endpoint of $T$. If the endpoint is
$A$, take $f=BC$ and $h=AB$; if it is $C$, take $f=AB$ and $h=BC$.
Again $e$ is disjoint from $f$, and the same inequality proves (i).
These cases exhaust the triple form.

Suppose instead that the only nonsingleton cells of $P$ are disjoint
adjacent pairs $f$ and $h$, with $F=D(f)+D(h)$. If $e$ equals either of
them, the chosen fine partition refines $P$. If $e$ is disjoint from at
least one of them, say $f$, then the partition with nonsingleton cells
$e,f$ is feasible at capacity $N-2$, and
\[
 D(e)+D(f)\le D(h)+D(f)=F.
\]
It is therefore a compatible flat optimum. Notice that this argument
allows $e$ to meet $h$, or to be disjoint from both $f$ and $h$.

The sole case remaining is that $e$ meets both $f$ and $h$. Write the two
coarse pairs in order as $AB$ and $CD$. An adjacent pair meeting both
must be $BC$. Thus $A,B,C,D$ are four consecutive original states,
and $P$ and $e$ have the form in (ii).

It remains to prove the exact flat-value identities and the price
comparison. The restricted quartet has a three-cell partition with pair
$BC$, of cost $q$, and the two-cell partition $AB\mid CD$, of cost $F$.
Conversely, any partition of $S$ with at most three (respectively two)
cells extends to a partition of the entire source with at most $N-1$
(respectively $N-2$) cells by making every state outside $S$ a singleton.
Its distortion is unchanged because every outside singleton has zero
cost. This rules out restricted values smaller than $q$ or $F$ and proves
the two identities. This step applies to arbitrary restricted partitions,
including noncontiguous ones.

Finally, apply the same singleton extension simultaneously to both levels
of any restricted hierarchy. Refinement is preserved, distortion is
unchanged, and the denominators are precisely the same. The full-source
optimum can only be better than this feasible hierarchy, proving
$R_N\le R_4(S)$. Restricted optimal hierarchies exist because their
partition set is finite; equivalently one can take the infimum of the
same comparison over all of them. For $N=3$ the coarse capacity is one,
whose single cell is compatible with every fine partition.
\end{proof}

\begin{corollary}[Equality of the all-size and quartet envelopes]
\label{cor:envelope}
For a fixed $C^2$ strictly convex generator $\varphi$ on $I$, let
$\mathcal C_N(\varphi)$ be the supremum of $R_N$ over $N$ distinct
locations in $I$ and positive weights. Then
\[
 \mathcal C_N(\varphi)\le\mathcal C_4(\varphi)\quad(N\ge3),
 \qquad
 \sup_{N\ge3}\mathcal C_N(\varphi)=\mathcal C_4(\varphi).
\]
The same equality holds after taking a supremum over any class of such
generators, in particular nonnegative nonzero polynomial curvatures of
degree at most $r$ on $[0,1]$.
\end{corollary}

\begin{proof}
Every hierarchy price is at least one, and the quartet envelope is at
least one. Apply Theorem~\ref{thm:reduction} pointwise. Taking the
supremum over $N$ includes $N=4$, so the reverse inequality for the last
display is automatic. If probabilities summing to one are required,
divide every weight in the restricted quartet by the same total mass;
the price is unchanged by this common scaling. No individual cell is
conditionally renormalized.
\end{proof}

\begin{remark}[Meaning and limits of ``exact'']
The reduction has no multiplicative loss and preserves both flat
denominators exactly before any common weight normalization. It asserts
domination by a restriction, not equality of an individual source price
with the selected quartet price. The lifted quartet hierarchies form
only a subset of the global feasible hierarchies. The corollary asserts
equality of the supremum over all finite source counts with the quartet
supremum; it does not prove equality separately for every fixed $N>4$ and fixed
generator. Section~\ref{sec:embedding} instead constructs special polynomial
witnesses with exact deterministic price preservation.
\end{remark}

